\documentclass[10pt, aspectratio=169]{beamer}

\usetheme{metropolis}
\definecolor{alertc}{HTML}{008cd9}
\definecolor{myDarkTeal}{HTML}{1c2d30}
\setbeamercolor{alerted text}{fg=alertc}
\setbeamercolor{normal text}{fg=myDarkTeal,bg=black!2}
\metroset{block=fill}
\usepackage{appendixnumberbeamer}

\usepackage{booktabs}
\usepackage[scale=2]{ccicons}
\usepackage{array}
\newcolumntype{L}[1]{>{\raggedright\arraybackslash}m{#1}}
\newcolumntype{C}[1]{>{\centering\arraybackslash}m{#1}}
\newcolumntype{R}[1]{>{\raggedleft\arraybackslash}m{#1}}

\usepackage{siunitx}

\usepackage{xspace}
\newcommand{\themename}{\textbf{\textsc{metropolis}}\xspace}

\usepackage{amsmath}
\def\RR{\mathbb{R}}
\def\Pr{\mathrm{Pr}}
\DeclareMathOperator*{\argmax}{arg\,max}
\DeclareMathOperator*{\argmin}{arg\,min}
\usepackage{algorithm}
\usepackage{algpseudocode}

\usepackage{natbib}
\bibliographystyle{abbrvnat}
\setcitestyle{numbers,square}
\usepackage{graphicx}
\usepackage{xcolor}
\usepackage{hyperref}

\definecolor{mDarkTeal}{HTML}{23373b}
\definecolor{mLightBrown}{HTML}{EB811B}
\definecolor{mDarkBrown}{HTML}{B85002}
\definecolor{mLightGreen}{HTML}{14B03D}
\setbeamertemplate{itemize/enumerate subbody begin}{\normalsize}

\title{Beating the International Prognostic Index for high-risk DLBCL patients}
\subtitle{Master thesis final report}
\date{July 9th, 2024}
\author{Lukas Geßl}
\institute{Chair of Statistical Bioinformatics, Regensburg University}

\begin{document}

\maketitle

\section{Recap: The goal of this thesis}

\begin{frame}{}
  \begin{itemize}
    \item MMML-Predict: develop a cost-efficient classifier 
    that filters DLBCL patients with progression-free survival $\leq 2$ years more reliably 
    than the International Prognostic Index for non-Hodgkin's lymphoma (IPI).
    % % \pause
    \item The IPI \cite{ipi93} is a simple risk score (0--5) based on five clinical features. 
    The cohorts IPI $\geq i, i = 0, 1, \ldots, 5$, lack precision ($< 50 \%$) or 
    are too small to be clinically relevant (prevalence $< 10 \%$).
    % \pause
    \item Our classifier should label at least 15\% of patients as high-risk with a precision
    of at least $\max(50\%, \text{precision of IPI} \geq 4)$.
    % \pause
    \item Unlike the IPI, the new classifier can incorporate the whole range of 
    modern features (like transcriptomic, genetic, clinical data, already-existent 
    signatures) measured at diagnosis and even dynamic features measured during 
    the treatment.
  \end{itemize}
\end{frame}

\begin{frame}
  MMML-Predict will enroll 300 DLBCL patients in a prospective trial.
  \begin{itemize}
    \item Data for the first 200 patients \textit{will} arrive here and will be 
      our sole foundation to train classifiers and finally submit a single one.
    \item A group in Leipzig will test the submitted classifier on the remaining 
      100 patients.
  \end{itemize}
  % \pause
  For this thesis, we also play by these rules, but on already existing data.
\end{frame}

\section{How to find and sell the best model}

\begin{frame}{A two-step approach}
  \begin{description}
    \item[Validation] Of those models we have trained, we want to find and choose
       the model that performs best on new data to the best of \alert{our} knowledge.
    \item[Testing] We need to demonstrate the performance of the chosen model to 
      \alert{outside} people on new, independent data.
  \end{description}

  % \pause
  To this end, we split the data $(X, y)$ into a train cohort 
  $(X_{\text{train}}, y_{\text{train}})$ (also for validation) and test 
  cohort $(X_{\text{test}}, y_{\text{test}})$ (no more repeated splitting).
\end{frame}

\begin{frame}{Validation}
  We start with a set of tuples of hyperparameters $H$, where every $h \in H$ 
  defines a model up to its parameters.

  For every hyperparameter tuple $h \in H$, we
  % \pause

  \begin{enumerate}
    \item fit the model to the train cohort 
    in a cross-validation, yielding a vector of cross-validated predictions 
    $\hat{y}_\text{train} = \text{cv}(h)$.
    % \pause
    \item We use the cross-validated predictions to calculate the 
      cross-validated error $\text{err}(y_{\text{train}}, \hat{y}_\text{train})$.
  \end{enumerate}

  We select the model $m^*$ with hyperparameter tuple
  \begin{align*}
    h^* = \argmin_{h \in H} \ \text{err}(y_\text{train}, \text{cv}(h)).
  \end{align*}
\end{frame}

\begin{frame}{Testing}
  We calculate $m^*$'s predictions $m^*(X_\text{test}) = \hat{y}_{\text{test}}$ 
  on the test cohort and estimate its performance on independent data via
  \begin{align*}
    \text{err}(y_{\text{test}}, \hat{y}_{\text{test}}).
  \end{align*}

  % \pause
  For our problem, we choose $\text{err}(y, \hat{y})$ as the minimum of the 
  negative precisions with a prevalence of at least 17\% (model output usually 
  needs thresholding).

  % \pause
  Strictly speaking, the threshold for the model output is another hyperparameter, 
  but it is a platform-dependent one \citep{transplatform17}. On a new data set, 
  one might take the 17\% quantile of the model output as the threshold.
\end{frame}

\section{Let's talk about $H$: candidate models}

\begin{frame}{Model-agnostic hyperparameters \ldots}
  \ldots apply for every model. In our case, they concern the predictor matrix 
  $X \in \RR^{n \times p}$ and the response vector $y \in \{ 0, 1 \}^n \cup 
  (\RR \times \{0, 1 \})^n$.

  \begin{itemize}
    \item We add all combinations of at most $n_{\text{combi}}$ discrete features 
      that are positive in a share of at least $s_{\text{min}}$ patients to $X$; e.g. 
      we add a column ``female and ABC-type tumor'' if at least 5\% of patients 
      have this property.
    % \pause
    \item For $T > 0$, we provide the fitting algorithm a modified reponse $y$, namely
      \begin{itemize}
        \item for the binary response, we set $y_i = 1$ if the patient's progression-free 
          survival is $< T$, $y_i = 0$ otherwise,
        \item for the Cox response, we censor all samples with time to event 
          exceeding $T$ at $T$.
      \end{itemize}
    % \pause
    \item A-priori feature selection: which features do we include in $X$ in the 
      first place?
  \end{itemize}
\end{frame}

\begin{frame}{The most model-specific hyperparameter: model class}
  At the core, our models consist of 
  \begin{itemize}
    \item Cox proportional-hazards,
    \item logistic regression and 
    \item ordinary linear (or Gauss) regression 
  \end{itemize}
  models \cite{zerosumR}, % \pause with a loss function optionally equipped with
  \begin{itemize}
    \item $\ell_1$ or $\ell_2$ regularization,
    \item the zero-sum constraint on a subset of features \cite{transplatform17},
    \item standardization of the predictor.
  \end{itemize}

  % \pause
  Moreover, we deploy random forests \cite{ranger-gh}.
\end{frame}

\begin{frame}{Nested models}
  Given some ``early'' models $f_i: \RR^p \to \RR, i = 1, \ldots, m$, we can nest 
  them into another, ``late'' model $f: \RR^m \to \RR$ and get a new model 
  $f \circ (f_1, \ldots, f_m)$.

  % \pause
  \begin{itemize}
    \item Often, the early models have been trained on another data set, so we 
      observe their output as features in our data set (like the Lamis 
      signature): such $f_i$ are merely projections onto a feature.
    \item If we need to fit some of the early models to our data, how can we 
      get reliable cross-validated predictions for $f$? See next slide.
  \end{itemize}
  % \pause

  Typically, we train the early  model on the high-dimensional 
  part of the data (like gene expression) and use its output together with the 
  remaining features as input for the late model.
\end{frame}

\begin{frame}{}
  \begin{algorithm}[H]
    \caption{Nested pseudo cross validation} \label{alg:nested-pcv}
    \begin{algorithmic}[1]
      \State \textbf{Input:} Predictor matrix $X$, response $y$, 
        hyperparameter tuple $h = (h_1, h_2)$
      \State Fit $f_1$ to $(X; y)$ subject to $h_1$ in a $k$-fold cross-validation, 
        yielding cross-validated predictions $\hat{y}^{(1)}$.
      \State Fit $f$ to $(\hat{y}^{(1)}, f_2(X), f_3(X), \ldots, f_m(X); y)$ 
        subject to $h_2$ in a $k$-fold cross-validation, yielding 
        cross-validated predictions $\hat{y}$.
      \State $g \gets f \circ (f_1, \ldots, f_n)$
      \State \textbf{Output:} $(\hat{y}, g)$
    \end{algorithmic}
  \end{algorithm}

The \text{pseudo} cross-validated prediction for every sample in $\hat{y}$ slightly 
depends on the sample itself. Benefit: save a factor $k$ in time complexity.

Procede greedily (first tune $h_1$, then $h_2$) to avoid 
overfitting of cross-validated predictions to the training cohort.
\end{frame}

% \section{From theory to practice: software}

\begin{frame}{The R package patroklos}
  \begin{minipage}{0.15\textwidth}
    \centering
    \includegraphics[width=\textwidth]{figs/logo.png}
  \end{minipage}
  \hfill
  \begin{minipage}{0.8\textwidth}
    \raggedright
    patroklos~\citep{patroklos} solves this and analogous problems with the 
    presented methods.
  \end{minipage}
\end{frame}

\section{How this plays out on real data}

\begin{frame}{Meet the data}
  \begin{table}[ht]
\small
\centering
\begin{tabular}{L{.18\textwidth}R{.2\textwidth}R{.28\textwidth}R{.18\textwidth}}
  \hline
  & \textbf{Schmitz \cite{schmitz18}} & \textbf{Reddy \cite{reddy17}} & \textbf{Lamis test \cite{staiger20}} \\
  \hline
  \textbf{\# samples} & 229 & 604 & 466 \\
  \textbf{\# genes} & 25066 & 13302 & 145 \\
  \textbf{technology} & RNA-seq & RNA-seq & NanoString \\
  \textbf{high risk [\%]} & 36.6 & 31.5\footnotemark & 24.3 \\
  \textbf{IPI-45 prev. [\%]} & 12.9 & 21.6 & 17.0 \\
  \textbf{IPI-45 prec. [\%]} & 65.2 & 54.1 & 38.2 \\
  \textbf{other features\footnotemark} & signature ``genetic subtype'' & 
    genetic events: high expression, translocation, mutation & LAMIS signature \\
  \hline
\end{tabular}
\end{table}

  \footnotetext[1]{High risk is defined as overall survival < 2.5 years.}
  \footnotetext[2]{All datasets include the IPI features in thresholded form, 
    gender, cell of origin, and the LAMIS signature.}
\end{frame}

\begin{frame}{Intra-trial: Validate and test on the same data set}
  \begin{table}
    \centering
    \begin{tabular}{lrrr}
        \hline
         & \textbf{Schmitz} & \textbf{Reddy} & \textbf{Staiger} \\
        \hline
        \textbf{\# samples} & \num{58} & \num{151} & \num{117} \\
        \textbf{high risk [\%]} & \num{37.0} & \num{31.6} & \num{24.3} \\
        \textbf{$\text{prev}(\text{tIPI})$, $\text{prec}(\text{tIPI})$ [\%]} & \num{17.0}, 
            \num{50.0} & \num{19.2}, \num{42.1} & \num{39.2}, \num{38.7} \\
            \textbf{$\text{prev}(\mtc{i})$, $\text{prec}(\mtc{i})$ [\%]} & \num{35.1}, \num{68.4} & 
            \num{23.0}, \num{55.6} & \num{34.6}, \num{45.9} \\
        \textbf{precision \num{95}\%-CI $\mtc{i}$ [\%]} & $[\num{43.5}, \num{87.4}]$ & 
            $[\num{35.3}, \num{74.5}]$ & $[\num{29.5}, \num{63.1}]$ \\
        \textbf{HR $\mtc{i}$} & \num{1.00} & \num{1.04} & \num{13.6} \\
        \textbf{HR \num{95}\%-CI $\mtc{i}$} & $[\num{1.00}, \num{1.00}]$ & 
            $[\num{1.02}, \num{1.06}]$ & $[\num{1.99}, \num{93.2}]$ \\
        \textbf{logrank p $\mtc{i}$} & \num[scientific-notation=true]{3.69e-4} & 
            \num[scientific-notation=true]{1.82e-3} & 
            \num[scientific-notation=true]{9.38e-4} \\
        \hline
    \end{tabular}
    \caption{Statistics of intra-trial experiments on the respective test cohort. $\mtc{i}$ denotes the 
        model with minimal validation error on the respective training cohort.}
    \label{table:intra-trial}
\end{table}
\end{frame}

\begin{frame}{$m^*$'s architecture in a nutshell}
  \small
  \begin{columns}[T]
    \column{.3\textwidth}
      \textbf{Schmitz}\par
      Nested model as in Alg. \ref{alg:nested-pcv} with
      \begin{itemize}
        \item the early model (Gauss) trained on the 
            RNA-seq features,
        \item the late model (Cox) trained on the early model's output plus the 
          remaining features (IPI in all verions), $n_\text{combi} = 2$.
      \end{itemize}
    \column{.3\textwidth}
      \textbf{Reddy}\par
      Nested model as in Alg. \ref{alg:nested-pcv} with
      \begin{itemize}
        \item the early model (Gauss) trained on the 
            RNA-seq features,
        \item the late model (Cox) trained on the early model's output plus the 
          remaining features (five IPI features discretized), 
          $n_\text{combi} = 3$.
      \end{itemize}
    \column{.3\textwidth}
      \textbf{Lamis test}\par
      A logistic model trained on all features except for the NanoString gene 
      counts, $n_\text{combi} = 1$.
  \end{columns} 
\end{frame}

\begin{frame}{$m^*$ seems to the winner of a lottery: Schmitz}
  \begin{figure}[h]
    \centering
    \includegraphics[width=.9\textwidth]{../../results/schmitz/meta_gauss-glm.jpeg}
  \end{figure}
\end{frame}

\begin{frame}{$m^*$ seems to the winner of a lottery: Reddy}
  \begin{figure}[h]
    \centering
    \includegraphics[width=.9\textwidth]{../../results/reddy/meta_gauss-glm.jpeg}
  \end{figure}
\end{frame}

\begin{frame}{$m^*$ seems to the winner of a more predictable lottery: Lamis test}
  \begin{figure}[h]
    \centering
    \includegraphics[width=.9\textwidth]{../../results/staiger/meta_gauss-glm.jpeg}
  \end{figure}
\end{frame}

\begin{frame}{Inter-trial: Train and validate on one data set, test on another}
  \begin{table}[ht]
    \small
    \centering
    \begin{tabular}{lrrr}
        \hline
        & \textbf{Schmitz} & \textbf{Reddy} & \textbf{Lamis test} \\
        \hline
        \textbf{Schmitz} & (\num{12.9}/\num{65.2}) & (\num{17.7}/\num{59.6}) & (\num{17.1}/\num{50.7}) \\
        \textbf{Reddy} & (\num{17.8}, \num{71.1}) & (\num{21.6}/\num{54.1}) & (\num{18.0}/\num{53.2}) \\
        \textbf{Lamis test} & (\num{17.4}/\num{75.7}) & (\num{22.5}/\num{50.4}) & (\num{17.0}/\num{38.2}) \\
        \hline
    \end{tabular}
    \caption{Rows $i$ hold training cohorts, columns $j$ hold test cohorts. Diagonal 
        entries $(i, i)$ hold (prevalence/precision) of $\text{IPI} \geq 4$
        on cohort $i$. Off-diagonal entries $(i, j)$ hold (prevalence/precision) 
        on cohort $j$ of the best validated model $m^*_i$ trained on cohort $i$.}
    \label{tab:inter_trial}
\end{table}
\end{frame}

\begin{frame}{A closer look at $m^*$ for Reddy $\to$ Lamis test}
  We train a logistic model with $\ell_1$ penalty and standardization of 
  the predictor, for $T = \num{2.3}$ and 
  $n_\text{combi} = 2$, providing as features
  \begin{itemize}
    \item LAMIS score,
    \item cell of origin,
    \item IPI group: low (0--1), intermediate (2--3), high (4--5),
    \item the five thresholded IPI features,
    \item gender.
  \end{itemize}
\end{frame}

\begin{frame}{A closer look at $m^*$ for Reddy $\to$ Lamis test}
  \begin{table}[ht]
    \centering
    \small
    \begin{tabular}{lr}
        \hline
        \textbf{Feature} & \textbf{Coefficient} \\
        \hline
        Lamis score & \num{0.531} \\
        IPI group = intermediate \& age > 60 & \num{0.144} \\
        IPI group = low & \num{-0.266} \\
        IPI group = low \& gender = male & \num{-0.695} \\
        IPI group = low \& cell of origin = unclassified & \num{-0.033} \\
        IPI group = low \& ann arbor stage > 2 & \num{1.596} \\
        gender = male \& \# extranodal sites > 1 & \num{0.456} \\
        ABC/GCB unclassified \& performance status > 1 & \num{0.236} \\
        age > 60 & \num{0.379} \\
        age > 60 \& LDH ratio > 1 & \num{0.252} \\
        age > 60 \& performance status > 1 & \num{0.042} \\
        LDH ratio > 1 & \num{0.130} \\
        performance status > 1 & \num{0.904} \\
        \hline
    \end{tabular}
    \caption{Features with non-zero coefficients of the logistic model 
        $m^*_\text{Reddy}$.}
\end{table}
\end{frame}

\begin{frame}{A strong link between validation and test error: Reddy $\to$ Lamis test}
  \begin{figure}[h]
    \centering
    \includegraphics[width=.9\textwidth]{../../results/all/reddy_vs_lamis.jpeg}
  \end{figure}
\end{frame}

\begin{frame}{A strong link between validation and test error: Lamis test $\to$ Schmitz}
  \begin{figure}[h]
    \centering
    \includegraphics[width=.9\textwidth]{../../results/all/lamis_vs_schmitz.jpeg}
  \end{figure}
\end{frame}

\section{Conclusions and discussion}
\begin{frame}{Take aways}
  We wanted to deliver a classifier that defines a high-risk group of DLBCL 
  patients which is larger and more precise than that defined by the IPI.
  \begin{itemize}
    \item In intra-trial experiments, we could deliver on this promise for three 
      data sets. Inter-trial experiments worked even better.
    % \pause
    \item While simple, $\ell_1$-penalized models predicting from 
      high-dimensional gene expression levels only sometimes already beat 
      the IPI, one usually needs to integrate more features.
    % \pause
    \item Integrating \textit{already-existent} transcriptomic and genetic 
      signatures and the IPI features into another model reliably beats the IPI.
    % \pause
    \item Transferring these models from one data set (and platform) to another
      works very well (especially Reddy $\to$ Lamis test). Apparently, the
      size of the data set matters most.
  \end{itemize}
\end{frame}

\begin{frame}{Discussion}
  \begin{itemize}
    \item \textbf{Validation:} Ensure a reliable link between validated and 
      tested performance. How? % \pause
      \begin{itemize}
        \item Validating a smaller $H$ (proceding more greedily, relying on 
          prior, general knowledge).
        \item A refined cross validation following \citep{nested-cv-hastie} to 
          estimate the generalization error more reliably.
        \item Is our choice of $\text{err}$ too unstable? ROC-AUC isn't any 
          more stable.
        \item More samples.
      \end{itemize}
    % \pause
    \item \textbf{Training:} Deploy other, more complex models in the 
      integration step like boosted trees or neural networks. Balance 
      classification problem via sample weights in loss function.
    % \pause
    \item For MMML-Predict: rather more samples, less features.
  \end{itemize}
\end{frame}

\begin{frame}[standout]
  \begin{minipage}{0.45\textwidth}
    \centering
    \includegraphics[width=0.7\textwidth]{figs/pumuckl.png}
  \end{minipage}
  \hfill
  \begin{minipage}{0.50\textwidth}
      \centering
      Thank you! \par Questions?
  \end{minipage}
\end{frame}

\appendix

\begin{frame}[allowframebreaks]{References}
  \small
  \bibliography{../lit.bib}
  The ultra cute Pumuckl is taken from 
  \url{https://irp-cdn.multiscreensite.com/08191d67/dms3rep/multi/Pumuckl_Rennend.png}.
\end{frame}

\end{document}